\section{Conclusiones}
\label{sec:conclusiones}

La programación dinámica en C++ es la única combinación probada que resuelve una grabación entera en tiempo interactivo con garantía de optimalidad, por lo que es la recomendación del grupo para la práctica. Fuerza bruta tarda 21,3 s en $n = 32$ con $k = n/2$ y cada pulso adicional duplica ese tiempo, así que solo sirve como referencia de corrección. Backtracking con ambas podas resuelve esa misma instancia en microsegundos, pero sobre audio real su tiempo depende de dónde caen las repeticiones del tema y llega a ser 45 veces el de la programación dinámica en un prefijo de 350 pulsos. La misma tabla de la programación dinámica en Python tarda unas 100 veces más (73,2 s contra 0,668 s en $n = 1000$) y queda como verificación cruzada.

Cada técnica aportó algo que la anterior no tenía. Fuerza bruta fijó el costo óptimo contra el que se contrastaron los otros dos algoritmos, con un crecimiento medido de 2,034 por pulso contra 2,035 de la cota. Backtracking mostró que la poda de factibilidad divide los nodos por un factor fijo (2,48 en $n = 28$) sin cambiar la base del crecimiento. La de optimalidad baja los 67 millones de nodos de $n = 28$ a entre 45 y 261 según la familia, a cambio de un tiempo que depende de los datos. Que dos prefijos con el mismo último pulso y la misma cantidad de pulsos tengan las mismas extensiones llevó a la tabla $M[t][i]$, cuyo modelo de tiempo se cumplió a entre 0,44 y 0,47 ns por operación en sintéticas y reales.

Sobre audio real el modelo hace lo que promete y también muestra sus límites. En las 33 corridas la selección óptima cuesta entre un 28\,\% y un 94\,\% menos que truncar y descarta el material en uno o dos tramos largos (19 corridas con un salto y 14 con dos) unidos donde el tema se repite. En cuatro corridas, sin embargo, uno de los dos saltos descarta apenas uno o dos pulsos porque el costo no cuenta los cortes. En los diálogos el croma y los MFCC eligen recortes distintos, con índice de Jaccard 0 entre los pulsos descartados al 70\,\% en dos de las cuatro escenas, así que decidir cuál corte suena mejor exige una prueba de escucha que no se hizo.

Las mejoras posibles salen de las limitaciones del modelo y de las implementaciones. Si solo interesa el costo, la tabla de predecesores sobra y la programación dinámica corre con dos filas en $O(n)$ de memoria, mientras que precalcular los costos ahorra a lo sumo un factor $d$. En backtracking, sumar al costo parcial el salto más barato entre los pulsos restantes, cuando todavía hay que descartar alguno, daría una cota inferior que poda antes. Penalizar cada salto con una constante o exigir un largo mínimo por tramo descartado eliminaría los cortes de un pulso con la misma recurrencia, porque la penalización se suma a $c(i,j)$ y el largo mínimo solo quita valores del rango de $i'$. Combinar croma con MFCC es lo que convendría probar junto con esa prueba de escucha.
